\section{Remote databases as relations}
\label{sec:external-relations}

\begingroup
\setlength{\emergencystretch}{3em}
Public sequence databases differ in identifiers, response formats, and
retrieval models. NCBI separates taxonomy, assembly, gene, and sequence
services~\cite{ncbi_datasets_v2_api}. UCSC provides assembly and interval
retrieval~\cite{ucsc_rest_api}. UniProt serves protein records and
annotations~\cite{uniprot_programmatic_access}, while EBI supports cross-domain
discovery~\cite{ebi_search_rest_api}.
\par
\endgroup

BioSeqDB defines six remote relation shapes in the \texttt{bioseq\_external}
schema, but the generated table definitions currently omit the required
relation-entity option. An installation must add that option before the custom
scan can execute them. They cover taxonomy, provider catalogs, genes,
locations, complete sequences, and bounded sequence chunks.

When configured, these relations have no local heap data. PostgreSQL replaces
a physical scan with a custom scan~\cite{postgresql17_custom_scan}. Every
qualification remains a local PostgreSQL check. Supported equality and array
predicates are also lowered into a provider-neutral query that constrains the
remote request. The scan selects a provider adapter and emits ordinary typed tuples. A
remote record can therefore participate in a projection, filter, or join
without a separate download script.

\begin{table}[htbp]
  \centering
  \small
  \caption{External relation groups and their resource contracts.}
  \label{tab:external-relations}
  \begin{tabularx}{\linewidth}{p{.25\linewidth}p{.31\linewidth}X}
    \toprule
    Relation group & Result & Resource behavior \\
    \midrule
    Taxonomy and catalog & Normalized metadata rows & Provider response is
    decoded before row production \\
    Gene and location & One row per nested item & PostgreSQL applies residual
    predicates after normalization \\
    Sequences & One packed sequence value & Memory may depend on record length \\
    Sequence chunks & Ordered packed values & Accession-based NCBI retrieval
    uses a bounded queue \\
    \bottomrule
  \end{tabularx}
\end{table}

A custom scan retains state between tuple requests, postpones network access
until execution, and reports the chosen provider through \texttt{EXPLAIN}.
Ordinary \texttt{EXPLAIN} performs no network request, while \texttt{EXPLAIN
ANALYZE} consumes the remote stream.

\section{Provider and execution boundaries}
\label{sec:external-provider-boundary}

Provider integration combines a capability descriptor, a generated typed HTTPS
client, and a handwritten adapter. The client handles validation, encoding,
size limits, and transport. The adapter maps provider fields to a canonical
query and message model, keeping provider-specific formats outside PostgreSQL
planning.

The adapters have different capabilities. NCBI covers metadata and nucleotide
retrieval, UCSC covers catalogs and intervals, UniProt supplies protein
metadata, and EBI supports discovery. Protein sequence materialization is not
available through the nucleotide-oriented sequence type. A request that no
selected provider can serve fails explicitly. In a multi-provider scan, an
adapter error can be suppressed when another adapter returns rows, so such a
result may be partial.

At scan initialization, provider predicates select adapters, while supported equality
and constant array predicates constrain remote requests. Array predicate
expansion is currently capped at 64 equality-query combinations. Excess
combinations are omitted with a warning, so callers that require complete
results must remain below this limit. At execution time, the scan opens a
provider session on first demand. Metadata queries can collect a bounded set of
owning messages. Large NCBI nucleotide queries instead use an incremental
FASTA source and a queue bounded by both bytes and item count. Only the
PostgreSQL backend thread creates \texttt{bioseq.seq} values and writes tuple
slots. Network workers exchange ordinary owning values and never call the
PostgreSQL API.

\begin{algorithm}[htbp]
  \caption{External relation scan}
  \label{alg:external-relation-scan}
  \begin{algorithmic}[1]
    \Require relation kind, predicates, requested columns
    \State retain all predicates locally and lower supported predicates remotely
    \State select the requested or default adapter
    \State open the provider session when the first row is requested
    \State \textbf{while} the executor requests another row \textbf{do}
      \State obtain the next provider message and convert required sequence values
      \State place typed fields in the tuple slot for residual SQL processing
    \State \textbf{end while}
    \State cancel and close the session on rescan, error, or query completion
  \end{algorithmic}
\end{algorithm}

The streaming transport verifies its HTTPS destination, restricts redirects,
and applies connection and progress time limits. The archive path also applies
explicit wire-size and decoded-size limits. Incremental NCBI FASTA retrieval
bounds queue and carry memory but does not impose a total response-size limit.
The NCBI streaming scan checks PostgreSQL cancellation between bounded waits.
Retry is allowed only before the first row is emitted, which prevents duplicated
visible chunks.

Whole-record retrieval still requires record-sized memory, and archive
retrieval requires bounded buffering before decoding. Requests are not currently
rate-coordinated across backends. Chunked retrieval is preferable for large
records.
